\documentclass[10pt,a4paper]{amsart}
\usepackage[margin=19mm]{geometry}
\usepackage{amsmath,amssymb,mathtools}
\usepackage[scr=rsfs,cal=euler]{mathalfa}
\usepackage{xfrac}
\usepackage[unicode,hidelinks,bookmarks=false]{hyperref}
\newcommand{\ie}{i.e.\ }
\newcommand{\cf}{cf.\ }
\newcommand{\resp}{respectively\ }
\newcommand{\Subsection}{\subsection}
\providecommand{\nobreakdash}{\nobreak\hbox{-}}
\setlength{\emergencystretch}{2em}
\multlinegap0pt
\newtheorem{prop}{Proposition}
\def\E{\mathcal{E}}
\newcommand{\I}{\mathrm{I}}
\newcommand{\II}{\mathrm{II}}
\newcommand{\IV}{\mathrm{IV}}
\newcommand{\Ino}{\operatorname{{Ino}}}
\def\P{\mathbf{P}}
\newcommand{\Z}{\mathbf{Z}}
\newcommand{\caZ}{{\mathcal Z}}
\title{Higher modularity of elliptic curves over function fields}
\author{Adam Logan and Jared Weinstein, with an interlude by Masato Kuwata}
\date{Source: arXiv:2211.11149v2 (CC BY 4.0)}
\begin{document}
\maketitle

\noindent\emph{Source and licence.} Higher modularity of elliptic curves over function fields. Version arXiv:2211.11149v2 (CC BY 4.0), distributed by arXiv under the Creative Commons Attribution 4.0 International licence, http://creativecommons.org/licenses/by/4.0/, as recorded in the arXiv licence metadata for this exact version and shown on the abstract page https://arxiv.org/abs/2211.11149v2. Full licence notice: \texttt{licenses/arxiv-2211.11149-CC-BY-4.0.txt}.\par\smallskip

\noindent\textit{Editorial scope.} Counted unit: the article's prop prop (source0.tex lines 2389--2390). The article's complete original proof of it is delivered with it (source0.tex lines 2392--2453, one \texttt{proof} environment of 62 lines). No mathematics is written, repaired or re-derived: the statement and the proof are the article's own lines, in the article's own notation, and no retained line is substituted or re-flowed.  Retained uncounted as the source-local context the proof needs: lines 2383--2386.  One declared adaptation: exactly 2 ancillary strings inside the retained spans are omitted (image-inclusion commands and, where the source repeats a label, the repeated label definitions) and no other line differs from the source; the figure environments, with their placement, captions and labels, are kept, and the package records the omitted strings in fidelity receipts bound to the affected bodies.  No retained line contains a citation, so the wrapper carries no bibliography and no bibliography entry.  The retained lines contain the article's own diagram code, which the wrapper typesets with the standard tikz packages; no image file is delivered and the package declares no asset dependency.  Editorial formatting only, in addition: an AMS \texttt{amsart} preamble that restates the article's own declarations and macros used by the retained lines, namely the environments \texttt{prop} on separate counters, together with the standard packages those lines need.  The retained environments print in this document as Proposition 1 in order of appearance, whereas the article prints the same items with the numbering of its own full document.  The complete native source of the article as distributed in the arXiv e-print for this exact version is bundled unchanged as \texttt{sources/133760/source0.tex}.
\begin{equation}
\label{EqEIV}
 y^2 = x^3 + 9 x^2 + 24 t x + 16 t^2. 
 \end{equation}

\begin{prop} There is a morphism $\caZ^2(\E)\to \Ino(\E^2)$ of degree 9 of varieties over $U^2$.
\end{prop}

\begin{proof} 
%Write $\E^2_{\eta}=E_1\times E_2$ as usual.  
%%Twisting $E_{\IV^{*}}$ by $-3$ makes the fiber at $t=0$ nonsplit and the one at $t=1$ split.
%%So, take three copies of $E_{\IV^{*}}^{(-3)}$:
%%$E_{i} : y^2 = x^3 - 27 x^2 + 216 t_{i} x -432 t^2$ ($i=1,2,3$).
%The Inose surface $\Ino(E_{1}\times E_{2})$ is given by 
%\begin{multline*}
%\Ino(E_{1}\times E_{2}):
%Y^2 = X^3 - 27 (8 t_{1} - 9)(8 t_{2} - 9) X 
%\\
%- 54\Bigl(32 t_{1}^3 (t_{1} - 1) T - (8 t_{2}^2 - 36 t_{2} + 27) (8 t_{1}^2 - 36 t_{1} + 27) + \frac{32 t_{2}^3 (t_{2} - 1))}{T}\Bigr)
%\end{multline*}
The surface $\caZ^2(\E)$ is obtained by substituting $t=s(s-t_1-t_2+1)/(s-t_1t_2)$ into \eqref{EqEIV}, giving
\begin{equation}
\label{eq:base-ch-IV}
y^{2} = x^{3} +9  (s -t_{1} t_{2})^2 x^{2} 
+ 24 s (s - \delta) (s - t_{1} t_{2})^3 x+16 s^2 (s - \delta)^2 (s - t_{1} t_{2})^4,
\end{equation}
It has singular fibers of type $\IV^\ast,\IV^\ast,\I_3,\I_3,\I_1,\I_1$ at $s=\infty,t_1t_2,0,t_1+t_2-1,t_1,t_2$, respectively, with Mordell-Weil group again $\Z/3\Z$.  Let $\sigma_0$ be the identity section and let $\sigma_1$ be one of the 3-torsion sections.   We display here the $\IV^\ast$ and $\I_3$ fibers along with $\sigma_0$ and $\sigma_1$:

\begin{center}

\end{center}



This configuration contains two $\II^\ast$ fibers, highlighted below:
\begin{center}

\end{center}


Therefore there exists a genus 1 fibration $\caZ^2(\E)\to \P^1\times U^2$ with two $\II^{*}$ fibers.  
To perform a $3$-neighbor step, we first find the tangent line at each point of the $3$-torsion section $\sigma_{1}=\bigl(0,4s(s - t_{1} - t_{2} + 1)(s - t_{1}t_{2})^2\bigr)$:
\[
y=3(s-t_{1}t_{2})x + 4s(s - t_{1} - t_{2} + 1)(s-t_{1}t_{2})^2.
\]
Using this, we find an elliptic parameter 
\begin{equation}
\label{paramIV}
w = -\frac{t_{2}^3(t_{2} - 1)}{8}\cdot
\frac{y - 3(s-t_{1}t_{2})x - 4s(s - t_{1} - t_{2} + 1)(s-t_{1}t_{2})^2}{s^2(s-t_{1}t_{2})^4},
\end{equation}
which leads to a cubic curve in $x'$ and $s$ with parameter~$w$:
\begin{equation}
\label{EqIVCubic}
t_{2}^6 (t_{2} - 1)^2 {x'}^3 + 12 t_{2}^3 (t_{2} - 1) w (s -t_{1} t_{2}) x'  + 8 t_{2}^3 (t_{2} - 1) w (s - t_{1} - t_{2} + 1) - 8 w^2 s (s - t_{1} t_{2})^2,
\end{equation}
where $x'=2s(s-t_{1}t_{2})^2 x$.

This fibration lacks a section, but it does have a multisection of degree~$3$.  Indeed, any of the uncolored vertices in the figure represents a curve which meets each fiber with multiplicity~$3$.  Therefore there exists a degree~$9$ map from $\caZ^2(\E)$ to the Jacobian $J$ of the fibration.  The Weierstrass equation for $J$ is\[
Y^2 = X^3 - 3  (8 t_{2} - 9) (8 t_{1} - 9)X 
+ 2 (32 t_{1}^3 (t_{1} - 1) w - (8 t_{2}^2 - 36 t_{2} + 27) (8 t_{1}^2 - 36 t_{1} + 27) + 32 t_{2}^3 (t_{2} - 1) w^{-1}),
\]
which coincides with the twist of $\Ino(\E^2)$ by~$-3$.%
%\footnote{Adam: There seem to be a couple of sign errors.  I believe that what Masato meant to write is
%$$Y^2 = X^3 - 3  (8 t_{2} - 9) (8 t_{1} - 9)X 
%+ 2 (32 t_{1}^3 (t_{1} - 1) w - (8 t_{2}^2 - 36 t_{2} + 27) (8 t_{1}^2 - 36 t_{1} + 27) + 32 t_{2}^3 (t_{2} - 1) w^{-1}).$$}
%He actually wrote Y^2 = X^3 + 3  (8 t_{2} - 9) (8 t_{1} - 9)X 
% + 2 (32 t_{1}^3 (t_{1} - 1) w + (8 t_{2}^2 - 36 t_{2} + 27) (8 t_{1}^2 - 36 t_{1} + 27) + 32 t_{2}^3 (t_{2} - 1) w^{-1})

\end{proof}

\end{document}
